% H.264/HEVC式混合编码器的主要模块，按两篇标准概述的编码流程重绘；宽约150mm，跨栏使用。
% 灰底模块组成编码器内部的解码环路，与独立解码器执行相同的计算；白底模块只在编码端出现。
\begin{tikzpicture}[x=1cm,y=1cm,font=\figurefont\small,
  enc/.style={draw=FigureBlue,fill=white,rounded corners=1.8pt,line width=0.8pt,
    minimum width=21mm,minimum height=8mm,align=center,inner sep=1.2mm},
  dec/.style={enc,draw=FigureStroke,fill=FigurePaper},
  sum/.style={circle,draw=FigureStroke,fill=white,line width=0.8pt,
    minimum size=5.5mm,inner sep=0pt},
  arr/.style={-{Stealth[length=1.8mm]},draw=FigureStroke,line width=0.8pt,rounded corners=1mm},
  side/.style={arr,dashed},
  note/.style={text=FigureMuted,font=\figurefont\footnotesize}]
  \node[enc] (in) at (0,3) {当前帧\\$\bm Y_t$};
  \node[sum] (sub) at (2.4,3) {$-$};
  \node[enc] (tr) at (4.6,3) {整数变换};
  \node[enc] (q) at (7.2,3) {量化\\（量化参数）};
  \node[enc] (ent) at (10.6,3) {熵编码\\CAVLC/CABAC};
  \node[enc,minimum width=14mm] (bs) at (14.0,3) {码流};
  \node[dec] (iq) at (7.2,1.5) {反量化\\反变换};
  \node[sum] (add) at (7.2,0.2) {$+$};
  \node[dec] (lf) at (9.8,0.2) {环路滤波\\去块、SAO};
  \node[dec] (buf) at (12.6,0.2) {已解码帧\\缓存};
  \node[enc] (me) at (12.6,1.7) {运动估计};
  \node[dec] (pred) at (2.4,-1.3) {帧内预测或\\运动补偿};
  \draw[arr] (in) -- (sub);
  \draw[arr] (sub) -- node[above,note] {残差} (tr);
  \draw[arr] (tr) -- (q);
  \draw[arr] (q) -- node[above,note] {系数} (ent);
  \draw[arr] (ent) -- (bs);
  \draw[arr] (q) -- (iq);
  \draw[arr] (iq) -- (add);
  \draw[arr] (add) -- (lf);
  \draw[arr] (lf) -- (buf);
  \draw[arr] (buf) -- (me);
  \draw[arr] (in.north) -- (0,4.1) -| (me.north);
  \draw[side] (me.west) -| node[pos=0.75,left,note] {运动矢量、模式} (ent.south);
  \draw[arr] (buf.south) |- node[pos=0.75,above,note] {参考帧} (pred.east);
  \draw[arr] (pred.north) -- node[left,note] {预测$\overline{\bm Y}_t$} (sub.south);
  \draw[arr] (2.4,0.2) -- (add.west);
  \draw[side] (me.east) -- (14.1,1.7) |- (2.4,-2.3) -- (pred.south);
  \node[note,anchor=west] at (3.0,-2.6) {灰底：编码器内部的解码环路，与解码器逐位一致};
\end{tikzpicture}
